\section{Introduction}
A prepared causal experiment determines which future tests can be predicted from a visible history. Its predictive quotient is a quotient by executable tests, not a formal parameter manifold. To turn this information into a risk statement one must also determine which reports are actually acquired, which computations can be performed causally, and which data persist between physical calls. These distinctions become inseparable when the observer may wait for an informative report but has only finitely many retained labels.

We ask for a joint acquisition--memory law for one fixed future audit. The difficulty is twofold. First, selection changes the acquired law: a rare favorable region is usable only after its acquisition cost has been paid. Second, an observer's memory includes its routing and stopping mechanism, not just its final prediction. A label may or may not serve both roles, depending on whether termination is a function of the retained state or an event emitted by a transition. The resulting difference is operational and cannot be removed by changing notation.

Our central theorem applies to experiments having two independently verifiable certificates. One admitted reference model supplies an acquired-count inequality against every legal competing policy. One common parameter-independent acquisition program balances the unnormalized report and task-loss measures on every surviving history and every model. The converse and construction therefore use different quantifiers; no interchange of minimization and maximization is involved. The acquired profile is
\[
 \Qm_m(b)=\inf_{|C|\le m}\inf_{\substack{0\le h\le b\\\int h\,d\nu=1}}
       \int\bigl[w+\operatorname{dist}(f,C)^2\bigr]h\,d\nu,
\]
where $\nu$ is an actual reference report law, $f$ its conditional audit mean and $w$ its conditional audit variance. The productive budget $b$ is derived from the physical cost inequality.

Theorem~\ref{thm:main} first states the checkpoint converse and common-program upper independently. Under matched effort and identity exports it then proves exact autonomous laws for both stopping interfaces. Theorem~\ref{thm:approx} replaces exact balance by conditional, state-dependent report and task defects, retaining the expected physical budget rather than silently inflating it. Proposition~\ref{prop:sharp-approx} proves that the attempt factor in its task-error term and the denominator in its waiting bound cannot be removed. These results require neither independent attempts nor an everywhere positive conditional variance.

Theorem~\ref{thm:directional} gives a separate nonintervening consequence. A current, arbitrarily evolving hidden unit vector is observed through a Gaussian sensor, and is audited only after commitment. For $d\ge2$ the risk is comparable to
\[
 (1+\log b)^{-1/2}+m^{-2/(d-1)}.
\]
Here $m$ is the available terminal width, not the total typed-state budget. The theorem also identifies the exact obstruction at $m\le d$, the fixed-width limiting risk, and the leading noise constant when the output width grows. Its converse holds against all legal stopping and memory rules, not merely threshold rules. The proof derives both the radial reliability and the directional covering geometry from the original Gaussian kernel. This application does not randomize the target and does not assume a balancing physical command.

The variational profile itself has an established ancestry. When $w=0$, writing $a=1/b$ and $g=h/b$ gives $0\le g\le1$, $\int g\,d\nu=a$, and the objective is the normalized distortion of a fractionally retained mass $a$. This is the population trimming formulation underlying trimmed $k$-means and impartial trimming \cite{CGM97,GG99}. With $w\ne0$ it is the corresponding marked objective. We do not claim the density cap or low-distortion selection principle as new. Classical reject-option prediction \cite{Chow,DHZ}, invariant decision procedures \cite{Luschgy}, and quantum randomization \cite{AMTW} supply related tools. Approximate rejection sampling from a fixed independent batch \cite{BP23,SPJ26} has a different budget and output objective: here physical attempts can be history dependent, the criterion is an audited decision loss, and a single finite-state program must work for every admitted model.

The contribution is the causal, fully accounted passage from these raw certificates to a stopped risk law, its role-aware and perturbative completion, and the multidimensional sensing consequence. The standard projection, compactness, likelihood and covering arguments are identified in the proofs. The theorem is a sufficient-class result; it does not assert that arbitrary positive kernels possess a balanced reference. The predictive quotient is realized only under its explicit measurability conditions in Appendix~\ref{sec:quotient}. Actual marked geometry and singular masses are treated before the two raw realizations; effective execution and typed causal morphisms are kept separate from ideal Borel existence.
